% ---
% titre: Priorité des opération
% theme: NO
% chapitre: Nombres relatifs
% sous_chapitre: Priorité des opérations
% niveau: [10e]
% type: EVAL
% tags: [c-priorite-operations,c-nombres-decimaux, p-question-TS]
% difficulte: 1
% corrige: true
% ---

\question
Effectue les calculs suivants \textbf{en notant toutes les étapes intermédiaires}.

\begin{parts}
\vspace{20pt}\part[1\half] $30-(5-8)\cdot(10-8)= $ 

\vspace{10pt}{\footnotesize\raggedleft Espace pour ta démarche et tes calculs\par}
	\begin{solutionorgrid}[150pt]
		\begin{align*}
		30-(5-8)\cdot(10-8)&=\\
		30-(-3)\cdot(10-8)&=\\
		30-(-3)\cdot2&=\\
		30-(-6)&=\\
		36&
		\end{align*}
		
	\hfill\textbf{1,5\,pt} (-0,5\,pt par erreur)
	\end{solutionorgrid}
	
\vspace{20pt}\part[1] $6,8 + 15 \cdot 0,3 + 2,1=$

\vspace{10pt}{\footnotesize\raggedleft Espace pour ta démarche et tes calculs\par}
	\begin{solutionorgrid}[150pt]
		\begin{align*}
		6,8 + 15 \cdot 0,3 + 2,1&=\\
		6,8 + 4,5 + 2,1&=\\
		13,4&
		\end{align*}
	
	\hfill\textbf{1\,pt} (-0,5\,pt par erreur)
	\end{solutionorgrid}

\vspace{20pt}\part[1] $220 + (45:5)^2 - 10=$

\vspace{10pt}{\footnotesize\raggedleft Espace pour ta démarche et tes calculs\par}
	\begin{solutionorgrid}[150pt]
		\begin{align*}
		220 + (45:5)^2 - 10&=\\
		220 + 9^2 - 10&=\\
		220 + 81 - 10&=\\
		301 - 10&=\\
		291&
		\end{align*}
	
	\hfill\textbf{1\,pt} (-0,5\,pt par erreur)
	\end{solutionorgrid}
	
\vspace{20pt}\part[1] $\sqrt{144} + 7\cdot(92-85)=$

\vspace{10pt}{\footnotesize\raggedleft Espace pour ta démarche et tes calculs\par}
	\begin{solutionorgrid}[150pt]
		\begin{align*}
		\sqrt{144} + 7\cdot7&=\\
		12 + 7\cdot7&=\\
		12 + 49&=\\
		61&
		\end{align*}
	
	\hfill\textbf{1\,pt} (-0,5\,pt par erreur)
	\end{solutionorgrid}
	
	
\vspace{20pt}\part[1] $(13-7)^2 + 4 \cdot 9 - 3=$

\vspace{10pt}{\footnotesize\raggedleft Espace pour ta démarche et tes calculs\par}
	\begin{solutionorgrid}[160pt]
		\begin{align*}
		6^2 + 4 \cdot 9 - 3&=\\
		36 + 4 \cdot 9 - 3&=\\
		36 + 36 - 3&=\\
		72 - 3&=\\
		69&
		\end{align*}
	
	\hfill\textbf{1\,pt} (-0,5\,pt par erreur)
	\end{solutionorgrid}
	
\vspace{20pt}\part[1] $60 \cdot 3 - 12  + \sqrt{16}=$

\vspace{10pt}{\footnotesize\raggedleft Espace pour ta démarche et tes calculs\par}
	\begin{solutionorgrid}[160pt]
		\begin{align*}
		60 \cdot 3 - 12  + \sqrt{16}&=\\
		60 \cdot 3 - 12  + 4&=\\
		180 - 12  + 4&=\\
		168 + 4&=\\
		172&
		\end{align*}
	
	\hfill\textbf{1\,pt} (-0,5\,pt par erreur)
	\end{solutionorgrid}
	
\vspace{20pt}\part[2] $\left[2 - 5 \cdot (-2) + (-9)\right]^2 : 9=$

\vspace{10pt}{\footnotesize\raggedleft Espace pour ta démarche et tes calculs\par}
	\begin{solutionorgrid}[200pt]
		\begin{align*}
		\left[2 - 5 \cdot (-2) + (-9)\right]^2 : 9&=\\
		\left[2 - (-10) + (-9)\right]^2 : 9&=\\
		\left(2 + 10 - 9\right)^2 : 9&=\\
		3^2 : 9&=\\
		9 : 9&=\\
		1&
		\end{align*}
	
	\hfill\textbf{2\,pt} (-0,5\,pt par erreur)
	\end{solutionorgrid}
	
\vspace{20pt}\part[2] $(-2)^3 \cdot 5 - \left[16 - (-9):3\right]=$

\vspace{10pt}{\footnotesize\raggedleft Espace pour ta démarche et tes calculs\par}
	\begin{solutionorgrid}[200pt]
		\begin{align*}
		(-2)^3 \cdot 5 - \left[16 - (-9):3\right]&=\\
		(-2)^3 \cdot 5 - \left[16 - (-3)\right]&=\\
		(-2)^3 \cdot 5 - 19&=\\
		-8 \cdot 5 - 19&=\\
		-40 - 19&=\\
		-59&
		\end{align*}
	
	\hfill\textbf{2\,pt} (-0,5\,pt par erreur)
	\end{solutionorgrid}
\end{parts}
